\documentclass[a4paper,12pt]{article}

\usepackage[T2A]{fontenc}
\usepackage[utf8]{inputenc}
\usepackage[russian]{babel}
\usepackage{geometry}
\geometry{left=30mm, right=15mm, top=20mm, bottom=20mm}
\usepackage{setspace}
\onehalfspacing 

\usepackage{titlesec}
\usepackage{tocloft}
\usepackage{xcolor}
\usepackage[
    colorlinks=true, 
    linkcolor=blue, 
    urlcolor=blue, 
    citecolor=blue
]{hyperref}

\begin{document}


\begin{titlepage}
    \begin{center}
        Санкт-Петербургский государственный университет\\
        Искусственный интеллект и наука о данных\\
        
        \vfill
    
        {\fontsize{14}{16}\selectfont \textbf{Бабинская Арина Игоревна} \par}
        
        \vspace{1cm}
        
        {\fontsize{16}{18}\selectfont \textbf{Разработка и исследование методов анализа сложности кода на языке C++ с реализацией программного инструмента} \par}
        
        \vspace{0.5cm}
        {\fontsize{14}{16}\selectfont Отчёт о прохождении \\ Учебной (ознакомительной) практики \par}
        
        \vfill
        
        \begin{flushright}
            Научный руководитель:\\
            Старший преподаватель кафедры \\
            системного программирования \\
            Андреев Юрий Александрович
        \end{flushright}
        
        \vfill
        
        Санкт-Петербург\\
        2026
    \end{center}
\end{titlepage}

\newpage


\tableofcontents
\newpage


\phantomsection
\section*{Введение}
\addcontentsline{toc}{section}{Введение}

Современная разработка программного обеспечения предъявляет высокие требования к качеству и поддерживаемости кода. Одним из ключевых факторов, влияющих на эти характеристики, является сложность программного кода. Для её оценки разработано множество метрик и инструментов, позволяющих количественно анализировать структуру программ. В частности, для языка C++ существуют такие широко используемые решения, как Klocwork, SonarQube, Clang-Tidy и Lizard. Эти инструменты предоставляют различные подходы к анализу сложности, включая вычисление цикломатической сложности, оценку размера кода и выявление потенциально проблемных участков. Однако, несмотря на их распространённость, они в основном ориентированы либо на безопасность и общее качество кода, либо используют ограниченный набор метрик, не всегда полностью учитывающий специфику языка C++.

В то же время язык C++ обладает рядом особенностей, существенно влияющих на сложность программ: шаблоны, перегрузка функций и операторов, лямбда-выражения и другие современные конструкции. Существующие инструменты либо частично учитывают эти аспекты, либо вовсе игнорируют их, что приводит к неполной или искажённой оценке сложности. Таким образом, наблюдается разрыв между возможностями современных средств анализа и реальными характеристиками кода на C++, что затрудняет получение объективной картины его сложности и снижает эффективность принимаемых на основе анализа решений.

В рамках данной работы ставится задача исследования и разработки методов анализа сложности программ на языке C++, направленных на устранение выявленного разрыва. Основной целью является создание программного инструмента, способного учитывать специфические особенности языка и предоставлять расширенный набор метрик для более точной оценки сложности. Для достижения этой цели необходимо провести обзор существующих решений, систематизировать метрики сложности, выявить влияние конструкций языка C++ на сложность кода, а также сформулировать требования к новому инструменту и разработать соответствующие методы анализа.

Предлагаемый подход основан на анализе исходного кода с использованием абстрактного синтаксического дерева, что позволяет учитывать структурные особенности программ. В рамках работы планируется разработка архитектуры инструмента и реализация его прототипа на языке Python. Инструмент будет выполнять выделение ключевых конструкций программы, рассчитывать набор интегральных метрик сложности для функций и файлов, а также формировать отчёт в формате JSON. Полученные результаты будут сопоставлены с результатами существующих инструментов, что позволит оценить эффективность предложенного решения и определить направления его дальнейшего развития.


\newpage


\section{Постановка задачи}
Целью работы является разработка и исследование методов анализа сложности программ на языке C++, учитывающих особенности языка и обеспечивающих более точную и информативную оценку сложности кода.

\subsection*{Задачи:}
\begin{enumerate}
    \item выполнить обзор и сравнительный анализ существующих инструментов и метрик оценки сложности программ на языке C++, с выявлением их ограничений;
    \item сформировать и формализовать расширенный набор метрик сложности, учитывающий особенности языка C++ (шаблоны, перегрузку, лямбда-выражения и др.);
    \item спроектировать архитектуру программного инструмента анализа сложности и реализовать его прототип с использованием языка Python;
    \item провести экспериментальное исследование разработанного инструмента на наборе программ на языке C++ и выполнить сравнительный анализ с существующими решениями для оценки качества предложенного подхода.
\end{enumerate}
\end{document}
